\input{/home/mbenjamin/Documents/Lycee/LateX/head.tex}

% Mise en page
\usepackage[left=1.5cm, right=2.0cm, top=1.5cm, bottom=1.5cm]{geometry}

% Numerotation des pages / la derniere
\makeatletter
\renewcommand{\@evenfoot}%
        {}%\thepage \slash \pageref{LastPage}}
\renewcommand{\@oddfoot}{\@evenfoot}
\makeatother

% En-tête
\renewcommand{\chapitre}{Statistiques à 2 variables }
\renewcommand{\seance}{Introduction}
\renewcommand{\classe}{CIEL 2}



\begin{document}

{\bf Nom, Prénom : \makebox[0.5\linewidth]{\dotfill} \hfill CIEL2} \\

{\centering \color{red} \bf \large TP Stats2 -- Intro \\}


\vspace*{0.5cm}

On donne ci-dessous le nombre d'étudiants diplômés d'un Master en mathématiques et statistiques aux États-Unis$^1$ ainsi que le nombre de programmeurs informatiques en Californie$^2$ entre 2012 et 2021.

\smallskip

\hfill {\footnotesize $^1$ étude du Centre National pour les Statistiques sur l'Éducation (NCES) des États-Unis.}

\hfill {\footnotesize $^2$ étude du Bureau des Statiques sur l'Emploi (BLS) des États-Unis.}

\begin{center}
\begin{tblr}{
    colspec={lXXXXXXXXXX},
    stretch=1,
    cells={valign=m},
    vlines,hlines,
    row{1} = {bg = darkWhite},
    column{1} = {bg = darkWhite}
}
	Année & 2012 & 2013 & 2014 & 2015 & 2016 & 2017 & 2018 & 2019 & 2020 & 2021 \\
	Diplômés & 6246 & 6957 & 7273 & 7589 & 8451 & 9082 & 10443 & 11382 & 12041 & 12583 \\
	Programmeurs & 41540 & 38750 & 37990 & 38650 & 	37570 & 34050 & 29740 & 24400 & 21800 & 23960
\end{tblr}
\end{center}

\bigskip

On va représenter ces données sur Géogebra à l'aide des manipulations suivantes :

\begin{enumerate}
	\item Dans le champs de saisie, taper : \texttt{s1:=\{2012,2013,\dots\}} en complétant les pointillés jusqu'à 2021.
	
	On dispose maintenant d'une première série \texttt{s1}.
	\item On va maintenant saisir les autres données d'une autre manière.
	
	Dans \texttt{Affichage}, choisir \texttt{Tableur}
	
	\begin{enumerate}
		\item Dans une première colonne (vous pouvez mettre le titre sur la première case), entrer les nombre de diplomés un à un.
		
		\item Dans une deuxième colonne, faire de même pour le nombre de programmeurs.
		
		\item Selectionner toutes les valeurs de la série diplomés puis : Clic droit > \texttt{Créer} > \texttt{Liste}
		\item Faire de même avec la série des programmeurs.
	\end{enumerate}
	\item Pour créer un nuage de points :
	
	\begin{enumerate}
		\item Dans le champs de saisie taper : \texttt{nuage = (s1,l1)}
		
		Vous venez de placer des points correspondants au nombre de dplomés en fonction de l'année (recadrer si besoin pour les visualiser)
		\item Dans le tableau, selectionner l'ensemble des deux valeurs présentes dans les deux colonnes, puis : Clic droit > \texttt{Créer} > \texttt{Liste de points}
		
		{\it Avec cette méthode, tous les points ont obtenu un nom.}
	\end{enumerate}
	\item On souhaite connaître la moyenne des séries.
	
	\begin{enumerate}
		\item Saisir \texttt{moyannee = moy(s1)}
		\item De même, calculer la moyenne du nombre de diplomés.
		\item Placer le point dont l'abscisse est la moyenne des années et l'ordonnée la moyenne du nombre de diplomés.
		\item Selectionner l'ensemble des valeurs de la première colonne, puis dans le menu, choisir \texttt{Statistiques à une variable.}
		
		Plusieurs onglets apparaissent alors. Tester toutes les possibilités.
		\item Selectionner l'ensemble des valeurs des deux colonnes, puis dans le menu, choisir \texttt{Statistiques à deux variables.}
		
		Plusieurs onglets apparaissent alors. Tester toutes les possibilités.
	\end{enumerate}
\end{enumerate}


\end{document}

